\section{Fuzzy-Enhanced Risk Management and Action Smoothing}
To address the limitations of the rigid, step-function risk mechanisms implemented 
in our initial environment design—specifically the hard turbulence threshold that 
aggressively suppressed buying actions—this framework integrates a Mamdani-style 
Fuzzy Inference System (FIS). By transforming our previously defined crisp thresholds 
into smooth, continuous membership functions, the agent adapts fluidly to volatile 
regimes without suffering from jarring policy shifts. The fuzzy framework operates 
through two independent components: an Action Smoother and a Reward Shaper.

\subsection{Fuzzy Action Smoothing}
In our previous stock trading environment setup, a binary switch was utilized to 
filter market actions when turbulence $\tau_t$ crossed an arbitrary threshold 
$\tau_{thr}$, forcing the mathematical constraint $a_{t}^{\prime} = \min(a_{t}, 0)$. 
As observed during our baseline training runs, this constraint created severe 
optimization cliffs. The Fuzzy Action Smoother directly replaces this previous setup 
by mapping the raw action component $a_i \in [-1, 1]$ generated by the policy network
and the current market turbulence $\tau_t$ into a risk-adjusted, continuous action 
$a_{t, safe}$.

The inputs are fuzzified using overlapping triangular and sigmoidal membership functions, 
which are visualized in detail in Figure~\ref{fig:action_smoother_mfs}. The input spaces 
are partitioned as follows:
\begin{itemize}
    \item \textbf{Raw Action ($a_i$):} Partitions the continuous action spectrum into 
    linguistic states: $\mathcal{A} = \{\text{Sell}, \text{Hold}, \text{Buy}\}$.
    \item \textbf{Market Turbulence ($\tau_t$):} Evaluates global instability relative to 
    historical maximum boundaries: $\mathcal{T} = \{\text{Low}, \text{Medium}, \text{High}\}$.
\end{itemize}

\begin{figure}[htbp]
    \centering
    \includegraphics[width=0.95\textwidth]{figures/action_smoother_mfs.png}
    \caption{Membership functions for the Fuzzy Action Smoother, mapping the partitions 
    for the raw agent action input $a_i$ (left) and the current market turbulence index 
    $\tau_t$ (right).}
    \label{fig:action_smoother_mfs}
\end{figure}

The linguistic rule-base is defined to prioritize asset preservation over aggressive 
execution during market distress, replacing the previous hard clipping logic:
\begin{align*}
    \mathbf{R}_1: & \text{ IF } \tau_t \text{ is High } \Longrightarrow a_{t, safe} \text{ is Sell} \\
    \mathbf{R}_2: & \text{ IF } \tau_t \text{ is Medium } \wedge a_i \text{ is Buy } \Longrightarrow a_{t, safe} \text{ is Hold} \\
    \mathbf{R}_3: & \text{ IF } \tau_t \text{ is Medium } \wedge a_i \text{ is Sell } \Longrightarrow a_{t, safe} \text{ is Sell} \\
    \mathbf{R}_4: & \text{ IF } \tau_t \text{ is Low } \Longrightarrow a_{t, safe} \text{ acts as } a_i
\end{align*}

Defuzzification is performed online via the centroid method to yield a smooth, risk-mitigated 
scalar execution value.

\subsection{Fuzzy Reward Shaping}
Our original reward formulation relied on a linear volatility penalty term 
$-\lambda e_t v_t$. As identified in our initial analysis, this rigid subtraction often 
obscured positive trading rewards during macro-driven market regimes, distorting the gradient 
path and penalizing the agent even when defensive actions were correctly taken. The Fuzzy 
Reward Shaper revises this approach by mapping asset exposure $e_t$ and the Market Volatility 
Index ($v_t$) into a dynamic, non-linear risk deduction scalar $\phi_t \in [0, 1]$.

The input variables utilize specialized state partitions illustrated in 
Figure~\ref{fig:reward_shaper_mfs}:
\begin{itemize}
    \item \textbf{Portfolio Exposure ($e_t$):} Divided into $\{\text{Low}, \text{High}\}$ 
    states to measure total capital at risk.
    \item \textbf{Volatility Index ($v_t$):} Evaluates market fear through three human-centric 
    configurations: $\{\text{Calm}, \text{Stressed}, \text{Panic}\}$.
\end{itemize}

\begin{figure}[htbp]
    \centering
    \includegraphics[width=0.95\textwidth]{figures/reward_shaper_mfs.png}
    \caption{Membership functions for the Fuzzy Reward Shaper, showing the continuous state intervals defined for portfolio exposure $e_t$ (left) and the VIX volatility metric $v_t$ (right).}
    \label{fig:reward_shaper_mfs}
\end{figure}

The internal rule-base penalizes high capital concentration only when the macro environment 
justifies defensive actions:
\begin{align*}
    \mathbf{R}_5: & \text{ IF } v_t \text{ is Panic } \wedge e_t \text{ is High } \Longrightarrow \phi_t \text{ is Severe} \\
    \mathbf{R}_6: & \text{ IF } v_t \text{ is Stressed } \wedge e_t \text{ is High } \Longrightarrow \phi_t \text{ is Medium} \\
    \mathbf{R}_7: & \text{ IF } v_t \text{ is Calm } \vee e_t \text{ is Low } \Longrightarrow \phi_t \text{ is None}
\end{align*}

The final defuzzified penalty weight modifies the environment's internal feedback path cleanly, 
updating our previous reward structure to:
\[
    r_t = \log\left(\frac{V_{t+1} + \epsilon}{V_{t} + \epsilon}\right) - (\phi_t \cdot \lambda_{scale})
\]

In summary, transitioning from our baseline architectural constraints to an integrated 
fuzzy logic pipeline shifts the operational paradigm from rigid, non-linear safety cutoffs 
to continuous, context-aware risk stabilization. Rather than imposing blind, invariant 
structural or environmental penalties that obscure underlying policy performance, this 
approach dynamically balances risk exposure against current macro volatility, smooths 
extreme action profiles, and preserves critical gradient pathways across chaotic market 
conditions.